
\section{Minimizar o custo da configuração de uma matriz dado restrições de linhas e colunas}
%https://codeforces.com/gym/105677/problem/J
Considere uma matriz $N\times M$ composta por células pretas e brancas em que cada célula preta é representada pelo número $0$ e cada célula branca por um número no intervalo fechado de $[1,9]$. Temos dois tipos de restrições, a restrição horizontal e vertical, as duas restrições são aplicáveis somente para as células pretas. A restrição horizontal é determinada pela tripla $(x,y,S)$ tal que $(x,y)$ representa a célula preta em que a restrição é aplicada e $S$ indica que as células brancas contíguas a direita de $(x,y)$ devem ter soma igual a $S$. A restrição horizontal é determinada também por uma tripla $(x,y,S)$ tal que $(x,y)$ representa a célula preta em que a restrição é aplicada e $S$ indica que as células brancas contíguas abaixo de $(x,y)$ devem ter soma igual a $S$. Dado uma matriz inicial e um conjunto de restrições, qual o custo mínimo para fazer com que a matriz obedeça a todas as restrições tal que o custo é igual a a diferença absoluta do valor de cada célula branca na matriz original pelo valor na nova matriz. Além disso, cada célula branca é coberta por exatamente uma restrição horizontal e uma vertical.

Podemos resolver esse problema em $\mathcal{O}(MCMF)$. Modelamos uma rede em que cada restrição horizontal e vertical são vértices, consideramos cada célula branca $W_i$ como uma aresta entre um par de restrições $R_i$ e $C_i$ (restrição linha e coluna respectivamente). Mas cada aresta vai ser considerada como dois vértices $W_i$, $W'_i$ ou $W_{in}$, $W_{out}$. Ligamos arestas do $Src$ para cada restrição $R_i$ com capacidade igual a $R_i.v'$ (valor da restrição modificada - vai ser esclarecido) e custo $0$, fazemos o mesmo tipo de aresta entre $Snk$ e $C_i$. Ligamos arestas com capacidade igual a infinito (esse infinito pode ser $8$) entre cada $R_i$ e $W_i$ correspondente, o análogo entre $W'_i$ e $C_i$. Além disso, ligamos duas arestas entre $W_i$ e $W'_i$, a primeira com capacidade $W_i.v-1$ e custo $-1$ e outra com capacidade $8 - (W_i.v-1)$ e custo $1$. Executamos o $MCMF$ e a resposta é possível somente se $\sum R_i.v' = \sum C_i.v' = Max\_Flow$ e o custo mínimo é $\sum (W_i.v -1) ~-~ Min\_Cost$.

Cada $W_i$ pode ser unicamente determinado por um par $(R_i,C_i)$ e $W_i$ contribui com o seu valor para as restrições horizontais e verticais. Assim, $\sum R_i.v = \sum W_i.v = \sum C_i.v$, então faz sentido pensar que a cada $1$ de valor de $W_i$, ele vai restringir em $1$ o valor do seu $R_i$ e $C_i$ correspondentes. A partir disso, podemos pensar em uma modelagem em que cada $W_i$ serve como uma aresta entre um $R_i$ e um $C_i$, porque passar fluxo em $W_i$ (passar valor em $W_i$) resulta em restringir o valor de $R_i$ e $C_i$. Por fim, consideramos uma modelagem em que cada $R_i$ tem uma aresta com $Src$ com capacidade igual ao seu valor, o análogo entre cada $C_i$ e $Snk$, além disso cada $W_i$ representa uma aresta que liga um $R_i$ e $C_i$. Existe uma consideração importante, parece correto dizer que a capacidade de $W_i$ é $9$, mas se considerarmos que é nove, temos a escolha de passar zero de fluxo por $W_i$ e essa escolha significa dizer que o valor de $W_i$ é zero, mas isso não pode ser possível, pois $W_i.v \in [1,9]$. Para contornar isso subtraímos $1$ de cada $W_i$, agora cada $R_i$ e $C_i$ vão ter um novo custo $c'$ tal que $c' = c - qtd_{W}$, ou seja, o novo custo de cada restrição é o custo original menos a quantidade de células brancas na restrição. Agora passar $0$ de fluxo representa colocar a célula com valor $1$, o intervalo $[1,9]$ mapeia para o intervalo $[0,8]$.

Essa modelagem serviria para verificar se existe alguma configuração que atende todas as restrições, mas não calcula o custo mínimo. Vamos aplicar um custo para cada aresta, as arestas do tipo $(Src,R_i)$ e $(C_i,Snk)$ vão ter custo $0$. Para cada aresta $W_i$, vamos criar dois vértices $W_i,W'_i$ e considerar arestas de $(R_i,W_i)$ e $(W'_i,C_i)$ com capacidade $8$ e custo $0$. Vamos também criar duas arestas de $W_i$ para $W'_i$, uma com capacidade $W_i.v'$ (subtrair $1$ de cada célula branca) e custo $-1$ e outra com capacidade $8-W_i.v'$ e custo $1$. A arestas principais são entre $W_i$ e $W'_i$, sempre que possível vamos passar o máximo de fluxo na aresta de custo $-1$ (queremos minimizar o custo, assim vamos querer saturar a aresta, ou seja, passar o valor original da célula). Se vamos passar mais fluxo que $W_i.v'$, temos que passar na aresta de custo $1$. Essa modelagem é genial porque para valores menores ou iguais a $W_i.v'$, sempre vamos querer passar o máximo de fluxo (porque o custo é negativo), mas para valores maiores que $W_i.v'$, vamos querer passar o mínimo de fluxo, assim o fluxo sempre vai convergir para passar o mais próximo possível de $W_i.v'$ e assim minimizando o custo. Mas por que essa modelagem vai dar o custo mínimo?

Considere uma configuração que atenda as restrições, nessa configuração vão existir células em que o novo valor é menor ou igual ao original (considere o valor original como $x$ e o novo como $x'$) e outras em que o novo valor é maior que o original (considere o valor original como $y$ e o novo com $y'$). Podemos dizer que $Custo = \sum(x - x') + \sum(y'-y)$, mas também podemos dizer que $Custo = \sum(x-x') + \sum(y'-y + (y-y)) = \sum(x-x') + \sum(y'-y) + \sum(y-y) = \sum(x+y) + \sum(y' - y) + \sum(-x')$. O termo $\sum(x+y)$ refere-se ao somatório do valor das células original que é constante, então minimizar $Custo$ equivale a minimizar $\sum(y' - y) + \sum(-x')$ e esse é exatamente a modelagem de custo da nossa rede. O termo $\sum(y'-y)$ equivale a saturar a aresta de custo negativo e passar uma quantidade de fluxo na aresta de custo $1$ e o termo $\sum(-x')$ é passar fluxo na aresta de custo $-1$. Assim a resposta vai ser $\sum W_i.v - Min\_Cost$.

\inputwithhash{./solutions/Problemas/RecoveringtheTablet.cpp}
